\section{Method}\label{sec:method}

We introduce an evaluation framework that measures the fidelity of
synthetic financial return series by decomposing the comparison between
real and synthetic returns into a set of metrics, each targeting a
distinct empirical regularity of asset returns \citep{cont2001empirical}.
The design of these metrics is guided by the desiderata of
\citet{raisa2025position} for synthetic-data evaluation measures.

The framework targets positive empirical regularities---properties
exhibited by real returns, such as heavy tails and volatility
clustering---rather than absence-based criteria such as zero linear
autocorrelation of raw returns.
Such criteria have weak discriminative power because even a structurally
inadequate generator such as geometric Brownian motion satisfies them by
construction.

The object of evaluation is a generator: any procedure producing
synthetic returns for a fixed cross-section of assets observed on a
common intraday clock.
The framework itself is generator-agnostic, and any generator satisfying
this input--output specification can be benchmarked.

Throughout the paper, a \emph{panel} is a matrix of log returns
$X \in \mathbb{R}^{T \times N}$ whose columns are $N$ assets and whose
rows are $T$ intraday timestamps on a common market clock; entry
$X_{t,i}$ is the log return of asset $i$ over minute $t$.
The framework takes as input real and synthetic panels aligned to the same intraday clock.
The panels need not share ticker identities, because real and synthetic ticker subsamples are drawn independently from their respective panels.
Both panels undergo the same preprocessing pipeline described in (Section~\ref{sec:experiments}).
For each metric, the framework asks whether the synthetic returns
reproduce the corresponding statistical property of the real returns.
Because two resamples of the observed real ticker corpus do not agree
exactly, their discrepancy provides an empirical reference under the chosen
sampling protocol.
The score compares the synthetic--real gap with this corpus-resampling gap (Section~\ref{sec:method_scoring}).

\subsection{Metric Architecture}\label{sec:method_architecture}

Every metric in the framework shares a common construction.
A metric $k$ consists of two components: a \emph{feature map} $\phi_k$,
which extracts from a panel of returns the statistical signature of a
stylized fact (for instance, the empirical quantile function of pooled
returns, or the autocorrelation function of absolute returns), and a
\emph{distance} $\rho_k$, which measures the discrepancy between two
such signatures.
The gap between two panels $\mathcal{A}$ and $\mathcal{B}$ under metric
$k$ is the pullback of $\rho_k$ through $\phi_k$:
\begin{equation}
  g_k(\mathcal{A}, \mathcal{B})
  = \rho_k\bigl(\phi_k(\mathcal{A}),\, \phi_k(\mathcal{B})\bigr).
  \label{eq:pullback}
\end{equation}

This construction has two practical consequences.
First, each gap is diagnostic: because $\phi_k$ targets a specific
stylized fact, a large gap identifies the distributional or temporal
structure responsible for the discrepancy rather than merely signalling
an unspecified difference.
Second, the framework is extensible: an additional property can be
incorporated by defining a new feature map and distance without modifying
the scoring or inference machinery of
Section~\ref{sec:method_scoring}.

Because $\rho_k$ compares only the extracted signatures, $g_k$ is blind
by construction to everything the feature map discards: two panels that
differ substantially can nonetheless satisfy
$g_k(\mathcal{A}, \mathcal{B}) = 0$ whenever they share the same
signature.
This blindness is deliberate.
Each feature map is intentionally insensitive to structures outside its
target.
For example, the marginal-distribution map is unchanged by reordering
returns, so temporal dependence is invisible to it.
No individual metric can therefore certify overall fidelity; the
framework instead assesses complementary properties separately.

\paragraph{M1: Tail-weighted marginal distribution}
Real returns exhibit heavy-tailed marginal distributions: extreme
returns occur far more frequently than under Gaussianity
\citep{cont2001empirical}.
M1 compares the unconditional return distributions of two panels, with
increased sensitivity to discrepancies in the tails.
The feature map first divides the returns of each ticker by their sample
standard deviation, computed with Bessel's correction over the finite returns of
that ticker in the panel; the ticker mean is not subtracted.
Without this step, pooling returns across tickers with unequal variances yields a
scale mixture that is leptokurtic even when every individual series is Gaussian,
so the pooled tail would partly reflect cross-sectional volatility dispersion
rather than distributional shape.
The map then pools the standardized returns within each panel and evaluates the
empirical quantile function $\hat{Q}(u)$ at $J$ equally spaced
probability levels
$u_j=(j+\tfrac{1}{2})/J$, $j=0,\ldots,J-1$.
This excludes $u=0$ and $u=1$, whose estimates correspond to the
sample minimum and maximum and are therefore the least stable order
statistics.
The distance is a tail-weighted $L1$ discrepancy between empirical
quantile functions, motivated by the quantile representation of the
one-dimensional Wasserstein-1 distance
\citep{villani2009optimal}:
\begin{equation}
g_1(\mathcal{A},\mathcal{B})
= \frac{1}{J}\sum_{j=0}^{J-1}
\widetilde{w}(u_j)
\left|
\hat{Q}_{\mathcal{A}}(u_j)
-
\hat{Q}_{\mathcal{B}}(u_j)
\right|
\label{eq:m1}
\end{equation}
The normalized weights are defined as
\begin{equation}
\widetilde{w}(u_j)
= \frac{w(u_j)}
{J^{-1}\sum_{k=0}^{J-1}w(u_k)},
\qquad
w(u)
= 1+\lambda
\left[
u^{-\alpha}+(1-u)^{-\alpha}
\right].
\label{eq:m1_weights}
\end{equation}
Here, $\alpha>0$ controls how strongly the weighting concentrates near
the distributional boundaries, while $\lambda>0$ controls the strength
of the tail component relative to the uniform component; configured
values are reported in Section~\ref{sec:experiments}.
Normalizing the weights to unit mean prevents the quantile-grid
resolution $J$ from mechanically changing the overall scale of the
distance; the weighting therefore redistributes sensitivity toward the
tails rather than rescaling the metric.

We use a smooth weighting rather than a hard tail indicator because a
hard cutoff is sensitive to small shifts of probability mass across the
cutoff under resampling and discards discrepancies in the bulk of the
distribution.
Because the unconditional distribution does not depend on temporal
ordering, any reordering of returns within a series leaves $g_1$
unchanged; temporal structure is therefore invisible to this metric by
construction, in the sense of
Section~\ref{sec:method_architecture}.

\paragraph{M2: Volatility clustering}
Real returns exhibit volatility clustering, whose quantitative signature
is the slow decay of the autocorrelation of absolute returns over long
lags \citep{ding1993long, cont2001empirical}.
We use absolute rather than squared returns because the autocorrelation
of $\lvert r_t\rvert^d$ is typically strongest and most persistent for
$d$ near one \citep{ding1993long}.
For each asset separately, the feature map computes the sample
autocorrelation of absolute returns at the selected lags and then
averages the resulting curves across assets; we write $\bar{a}(k)$ for
this cross-sectional average at lag $k$.
With 389 valid within-session returns, the largest admissible lag is
$K=388$, while the evaluated lag window is
$\mathcal{K}={k_{\min},\ldots,k_{\max}}$, with $k_{\max}\leq K$.

Two estimator details matter for intraday data.
First, the lag-$k$ products in the numerator are accumulated strictly
within trading sessions: observations separated by an overnight
boundary are not treated as lag-$k$ pairs, since including them would
contaminate the estimate with cross-session dependence, particularly at
long lags.
Second, the centring and variance normalization are computed from the
full sample of each asset, without correcting for the decreasing number
of valid pairs at larger lags.
This produces the intended shrinkage toward zero at long lags of the
biased-style autocorrelation estimator.

The distance is the average absolute difference between the
cross-sectionally averaged curves over the selected long-lag window:
\begin{equation}
g_2(\mathcal{A},\mathcal{B})
= \frac{1}{|\mathcal{K}|}
\sum_{k\in\mathcal{K}}
\left|
\bar{a}_{\mathcal{A}}(k)
-
\bar{a}_{\mathcal{B}}(k)
\right|,
\label{eq:m2}
\end{equation}
where the configured values of $k_{\min}$ and $k_{\max}$ are reported in
Section~\ref{sec:experiments}.
The window focuses on long lags because short-lag autocorrelation of
absolute returns can also be reproduced by short-memory volatility
models and is therefore less discriminative of persistent volatility
clustering.

\paragraph{M3: Aggregational Gaussianity}
The distribution of returns changes systematically under temporal
aggregation: at fine sampling scales returns are typically strongly
leptokurtic, while their distribution tends toward Gaussianity as returns
are summed over longer horizons \citep{cont2001empirical}.
The trajectory of this convergence contains information that a marginal
comparison at a single scale cannot capture: two processes may share a
one-period distribution yet aggregate differently because of differences
in their temporal dependence.

The feature map aggregates one-minute log returns into non-overlapping
$h$-minute blocks by summation, for each horizon $h$ in a set
$\mathcal{H}$ that includes the base scale $h=1$.
Aggregation is confined within trading sessions, so no aggregated return
spans an overnight boundary, and blocks containing missing observations
are discarded rather than partially summed.
\edited{All horizons are evaluated on the same contiguous within-session
support, so that the aggregation curve does not confound changes in scale
with changes in the time of day sampled.
For the horizons
$\mathcal{H}=\{1,5,15,30\}$, a session with 390 price observations provides
389 within-session return positions.
The largest common support divisible
by every horizon is therefore $30\lfloor389/30\rfloor=360$ positions.
We anchor this segment at the session close, retaining the final 360
return positions at every horizon.
This fixes the intraday interval used
for comparison and excludes the earlier return positions; the resulting
feature therefore describes aggregation over this closing segment rather
than the entire session.
Sessions shorter than the required support are
excluded from M3.
Within the retained support, blocks containing missing
observations remain excluded as described above.}

Aggregated returns are pooled across assets and sessions, and the bias-corrected sample
Pearson kurtosis $\hat{K}(h)=\hat{\kappa}(h)+3$ is computed at each scale, where $\hat{\kappa}(h)$ denotes sample excess kurtosis.

The feature is the normalized aggregation curve
$\hat{K}(h)/\hat{K}(1)$, which equals $1$ at the base scale and describes
the relative evolution of kurtosis across aggregation horizons.
Using Pearson rather than excess kurtosis avoids unstable normalization
for approximately Gaussian processes, whose base-scale excess kurtosis
may be close to zero.
The distance is the mean absolute difference between the two curves over
the scales at which both are defined:
\begin{equation}
g_3(\mathcal{A},\mathcal{B})
= \frac{1}{\lvert\mathcal{H}'\rvert}
\sum_{h\in\mathcal{H}'}
\left\lvert
\frac{\hat{K}_{\mathcal{A}}(h)}
{\hat{K}_{\mathcal{A}}(1)}
-
\frac{\hat{K}_{\mathcal{B}}(h)}
{\hat{K}_{\mathcal{B}}(1)}
\right\rvert,
\label{eq:m3}
\end{equation}
where $\mathcal{H}'\subseteq\mathcal{H}$ denotes the scales at which both
curves are finite.
Scales with too few valid aggregated observations are excluded rather
than estimated, with the observation floor and configured scale set
reported in Section~\ref{sec:experiments}.

Classical kurtosis is used because it corresponds directly to the
standard formulation of aggregational Gaussianity.
Because sample kurtosis remains sensitive to extreme observations, the
implementation also provides the quantile-based estimator of
\citet{moors1988} as an optional robustness variant, although it is not
used in the reported results.

\paragraph{M4: Regime-conditional tail shape}
The unconditional distribution compared by M1 averages over market
conditions.
For risk applications, however, it is also relevant whether the
heaviness of return tails varies across volatility regimes in a
realistic manner.
M4 therefore compares the conditional tail-shape profile across
relative volatility regimes.

The feature map proceeds in three steps.
First, for each asset, a strictly lagged volatility proxy
$\hat{\sigma}_{i,t-1}$ is computed as the rolling standard deviation of
prior returns over a trailing window.
The calculation is restarted at every session boundary, so that neither
the preceding session nor the current return $r_{i,t}$ can influence
the regime assigned to that return.
% >>> NEW (added for sensitivity analysis; defines min_frac, previously
% >>> in Table 1 but undefined in text). Please review wording and
% >>> verify the mechanism against the implementation --- the text
% >>> asserts exclusion, not imputation. Remove  when approved.
The trailing window is only partially populated near the start of
each session.
We retain a return only when the fraction of valid observations in its
window meets a minimum coverage threshold, and exclude it otherwise,
since a sparsely estimated proxy assigns regimes unreliably.
The proxy is used only to label observations; returns are not
standardized by it, since doing so would remove the conditional
variation that the metric is intended to capture.

Second, all valid volatility--return pairs are pooled across assets and
partitioned into a set $\mathcal{Q}$ of volatility-regime quantile bins,
with cut points computed separately from each panel.
This self-labeling evaluates whether real and synthetic data exhibit
similar tail behaviour across their respective relative volatility
regimes, rather than evaluating synthetic observations against
volatility levels defined by the real panel.

Third, within each regime $q$, a Generalized Pareto Distribution is
fitted by maximum likelihood, with location fixed at zero, to the
excesses $|r|-u_q$ conditional on $|r|>u_q$, where $u_q$ is a high
within-regime quantile \citep{mcneil2000estimation}.
The feature vector contains the estimated tail-shape parameters
$\hat{\xi}^{(q)}$ for $q\in\mathcal{Q}$.
Because the threshold is defined separately within each regime and only
the GPD shape parameter is retained, M4 compares conditional tail shape
rather than regime-specific exceedance frequencies or absolute tail
scales.
Absolute returns are used because the metric targets conditional tail
behaviour in magnitude; tail asymmetry is treated as a separate
property.

The distance is a weighted mean absolute difference between the two
tail-shape profiles,

\begin{equation}
g_4(\mathcal{A},\mathcal{B})
= \sum_{q\in\mathcal{Q}'}
\tilde{w}_q
\left|
\hat{\xi}_{\mathcal{A}}^{(q)}
-
\hat{\xi}_{\mathcal{B}}^{(q)}
\right|,
\qquad
\tilde{w}_q
= \frac{w_q}
{\sum_{j\in\mathcal{Q}'} w_j},
\label{eq:m4}
\end{equation}
where larger weights are assigned to higher-volatility regimes because
of their relevance for stress applications.
The weights are renormalized over the set
$\mathcal{Q}'\subseteq\mathcal{Q}$ of regimes for which both estimates
are defined.
A regime is excluded when it contains too few exceedances for stable
tail-shape estimation.
The exceedance floor, regime count, rolling-window length, minimum
rolling-window coverage, threshold quantile, and regime weights are
reported in Section~\ref{sec:experiments}.


\subsection{Scoring Protocol}\label{sec:method_scoring}

Raw gap values are difficult to interpret on their own because the metrics
have different units and numerical scales.
Moreover, independently estimated statistical signatures generally exhibit
nonzero discrepancies in finite samples.
We therefore compare each synthetic--real gap with an empirical real--real
baseline computed from two ticker resamples of the observed real corpus under
the same sampling protocol.
This baseline quantifies disagreement induced by resampling the available
corpus.

For each metric $k$, we estimate two quantities through repeated ticker
resampling.
An \emph{evaluation draw} is one round of resampling; we run $R$ draws,
indexed by $r=1,\ldots,R$.
In draw $r$, we sample $n$ assets with replacement twice from the real
panel, obtaining $\mathcal{R}_r^{(1)}$ and
$\mathcal{R}_r^{(2)}$, and once from the synthetic panel, obtaining
$\mathcal{S}_r$.
The sample size $n$ is fixed across draws and satisfies $n\leq N$.
Averaging over draws gives the \emph{real--real baseline gap} and the
\emph{synthetic--real gap},
\begin{equation}
  \bar{g}^{\,rr}_k
  = \frac{1}{R} \sum_{r=1}^{R}
    g_k\bigl(\mathcal{R}_r^{(1)}, \mathcal{R}_r^{(2)}\bigr),
  \qquad
  \bar{g}^{\,sr}_k
  = \frac{1}{R} \sum_{r=1}^{R}
    g_k\bigl(\mathcal{R}_r^{(1)}, \mathcal{S}_r\bigr).
  \label{eq:gaps}
\end{equation}
The first real sample $\mathcal{R}_r^{(1)}$ appears in both gaps, inducing correlation between the real–real and synthetic–real distances within each draw.
Sharing this reference sample reduces the Monte Carlo variability of their comparison. The $R$ matched draws provide Monte Carlo estimates of the average gaps induced by the resampling protocol, conditional on the observed real and synthetic panels.
Because sampling is performed with replacement, the two real samples may overlap.
These inner draws do not themselves quantify how the score would vary under a different ticker corpus; corpus-level uncertainty is estimated through the outer bootstrap described in Section~\ref{sec:method_ci}.

The similarity score for metric $k$ is
\begin{equation}
  s_k
  = \frac{\bar{g}^{\,rr}_k}{\bar{g}^{\,rr}_k + \bar{g}^{\,sr}_k}
  \in [0, 1].
  \label{eq:score}
\end{equation}
If both estimated gaps are exactly zero, the ratio is undefined and the
implementation returns $s_k=0.5$ by convention, reflecting equality of the
two estimated gaps.

The matched design---the same ticker-sample size $n$ in every
comparison---is essential.
Gap values fall as sample size grows, because larger samples estimate
each feature more precisely.
A baseline computed at a larger $n$ than the synthetic comparison would
therefore be too small, and the score would be biased downward; the
reverse holds for a smaller $n$.
Matching controls this sample-size effect by placing the real–real and synthetic–real gaps on the same finite-sample scale. It does not remove other sources of sampling or corpus variability.

The score is centred at $s_k=0.5$ in the algebraic sense that
$\bar g_k^{sr}=\bar g_k^{rr}$ if and only if $s_k=0.5$.
When $\bar g_k^{sr}>\bar g_k^{rr}$, the score falls below 0.5, indicating
that the synthetic--real discrepancy exceeds the empirical real-corpus
resampling gap.
When $\bar g_k^{sr}<\bar g_k^{rr}$, the score rises above 0.5, indicating
that the synthetic--real discrepancy is smaller than that resampling gap.

These comparisons are relative to the observed ticker corpus and the chosen
sampling protocol.
A value near 0.5 does not establish that the generator reproduces the
underlying real process or even that the extracted feature distributions are
equivalent.
Likewise, a value above 0.5 is not necessarily preferable: it may reflect
close agreement, but it may also arise from over-smoothing, reduced
cross-sectional variability, or other generator deficiencies.
The score should therefore be interpreted as a relative gap comparison,
rather than as a grade to be maximized, and should be considered together
with the underlying feature estimates and the remaining metrics.

\subsection{Outer-bootstrap confidence intervals}\label{sec:method_ci}

We quantify uncertainty with respect to ticker-corpus composition using an outer bootstrap. Let $X^R$ denote the preprocessed real panel with $N_R$ ticker columns, and let $X^{S,j}$ denote the corresponding preprocessed panel for generator $j$, with $N_{S,j}$ synthetic paths.

For each outer replicate $o=1,\ldots,O$, we sample $N_R$ real ticker columns with replacement from $X^R$ and independently sample $N_{S,j}$ synthetic columns with replacement from each $X^{S,j}$.
The same outer-resampled real panel is used for every metric and generator within a replicate.
On each resulting pair of outer-resampled corpora, we rerun the complete $R$-draw matched estimator of Section~\ref{sec:method_scoring} and obtain one score $s_k^{*(o)}$ for every metric–generator pair.

The reported 95\% confidence interval is the 2.5th–97.5th percentile range of the $O$ outer-bootstrap scores. The point estimate is computed separately on the original observed panels. Consequently, the point estimate is not mathematically required to lie inside the percentile interval.

The outer bootstrap is applied to the already preprocessed return panels. The resulting intervals therefore quantify sensitivity to ticker-corpus composition conditional on the observed time window, the preprocessing transformations estimated from the original panels, the fitted generators, and the available generated samples. They do not include uncertainty from refitting or retraining the generators.

Separately, we assess numerical stability by repeating the fixed-corpus matched estimator under independent inner-resampling seeds. The resulting standard deviation and range measure Monte Carlo sensitivity and are not interpreted as confidence intervals.

\subsection{Aggregate deviation score}\label{sec:method_aggregate}

The per-metric scores $s_1,\ldots,s_K$ are the framework's primary output,
because they identify the financial properties for which a generator exhibits
larger or smaller discrepancies.
For applications requiring a single summary, we additionally define an
optional aggregate deviation score.

For metric $k$, let
\[
    r_k = \frac{\bar g_k^{sr}}{\bar g_k^{rr}}
\]
denote the synthetic--real gap relative to the empirical real-corpus
resampling gap.
The ratio $r_k=1$ indicates equality of the two estimated gaps, whereas
$r_k>1$ and $r_k<1$ indicate deviations in opposite directions.
Because both directions represent departure from the reference value, we
define the symmetric metric-specific deviation
\[
    d_k
    =
    \exp\!\left(\left|\log r_k\right|\right)
    =
    \max\left(r_k,\frac{1}{r_k}\right)
    \geq 1.
\]

The aggregate deviation score is the geometric mean of these symmetric
deviations,
\[
    G_{\mathrm{dev}}
    =
    \left(\prod_{k=1}^{K}d_k\right)^{1/K}
    =
    \exp\!\left(
        \frac{1}{K}\sum_{k=1}^{K}\left|\log r_k\right|
    \right).
\]
It is computed in log space for numerical stability and is defined when all
metric-specific ratios are finite and strictly positive.

The value $G_{\mathrm{dev}}=1$ is obtained only when every included
metric has $r_k=1$.
Larger values indicate a greater average multiplicative deviation from the
empirical real-corpus reference, regardless of direction.
Unlike the geometric mean of the signed ratios, deviations above and below
one cannot cancel.

The aggregate remains secondary to the per-metric results.
It discards both the identity of the metric responsible for a discrepancy
and whether the corresponding ratio lies above or below one.
It must therefore be reported alongside the individual scores and gap ratios,
rather than as a standalone assessment of generator fidelity.

Uncertainty for $G_{\mathrm{dev}}$ is obtained directly from the outer
bootstrap.
Within each outer replicate, all metric-specific ratios are recomputed and
combined using the equation above.
The reported interval is the 2.5th--97.5th percentile range of the resulting
outer-bootstrap aggregate values; no additional bootstrap over the inner
evaluation draws is performed.
